\paragraph{Strategy.}
On the critical locus the twelve cut probabilities are linear in a $3\times3$ correlation matrix $T$, and
$\delta=\delta(T)$ is convex. The fidelity is the maximum of sixteen explicit plane functions. A critical state is
determined by $(r,s,w)$ with unit vectors $r,s$ (the $P=1$ coordinates) and a disc coordinate $w\in\mathbb R^2$,
$|w|^2=\tau$; the equality states sit at $w=0$. Analytic bounds settle $\tau\ge0.36$ (their crossing is at $\tau\approx0.35879$, TC-6); for $\tau\le0.37$ a
six-dimensional interval covering, together with a local cone lemma at the equality points, closes the proof.

\subsubsection*{Critical states as correlation matrices}
Let $M$ be the magic matrix, $v$ a unit critical normal vector, $A=\operatorname{Re}(vv^\dagger)$, and
$\psi=Mv$ the auxiliary two-qubit vector.  Put $T_{\alpha\beta}=\langle\psi|\sigma_\alpha\otimes\sigma_\beta|\psi\rangle$,
$t_{ji}=4\lambda_{ji}-1$ with $\lambda_{ji}$ the cut probabilities of Lemma~\ref{lem:fqe-critical-real-density}.
\begin{lemma}[Correlation-matrix model; analytic, exact check TC-1]\label{lem:fqe-critical-correlation}
\begin{enumerate}
\item The twelve rays $Mh_{ji}$ are Bell rays without local terms; their correlation matrices $C^{ji}$ are signed
permutation matrices (those of Lemma~\ref{lem:fqe-p1-cuts}), and $t_{ji}=\langle C^{ji},T\rangle$.
Consequently, by \eqref{eq:fqe-critical-real-tomography},
$A=\tfrac14I+\tfrac14\sum_{j,i}t_{ji}B_{ji}$: $A$ is an affine function of $T$.
\item $T=R_a\operatorname{diag}(\kappa,-\kappa,1)R_b^{T}$ with $R_a,R_b\in SO(3)$ and $\kappa=\sqrt\tau$.  Writing
$r=R_ae_3$, $s=R_be_3$ and $x_a,y_a$ ($x_b,y_b$) for the first two columns of $R_a$ ($R_b$),
\[
 T=rs^T+\kappa\bigl(x_ax_b^T-y_ay_b^T\bigr).
\]
Rotating the frames about $r$ and $s$ by $\theta_a,\theta_b$ turns $\kappa(x_ax_b^T-y_ay_b^T)$ into
$w_1(x_ax_b^T-y_ay_b^T)+w_2(x_ay_b^T+y_ax_b^T)$ with $w=\kappa(\cos(\theta_a+\theta_b),\sin(\theta_a+\theta_b))$.
\item For a plane $\Pi=VV^T$ of Lemma~\ref{lem:fqe-critical-real-density} the matrix $m_\Pi=V^TAV$ is affine in $T$ and
$\Phi=\max_\Pi\Phi_\Pi$, $\Phi_\Pi=\tfrac12\operatorname{tr}m_\Pi+\sqrt{\det m_\Pi}$.  Equivalently
$\Phi=\max\tfrac14\bigl(1+\sqrt{1-\tau}\,(\varepsilon\!\cdot\! r+\eta\!\cdot\! s)/\sqrt3+\varepsilon^TT\eta/3\bigr)$ over the
$32$ cube pairs of \eqref{eq:fqe-null-critical-bloch}.
\end{enumerate}
\end{lemma}
\begin{proof}
(1) The projector onto a maximally entangled ray is $\tfrac14(I+\sum C_{\alpha\beta}\sigma_\alpha\otimes\sigma_\beta)$,
with no local terms; the proof of Lemma~\ref{lem:fqe-p1-cuts} derives the twelve $C^{ji}$, and $t_{ji}=\langle C^{ji},T\rangle$
because $\lambda_{ji}$ is the expectation of that projector in $\psi$ (PM-1 and TC-1 recompute them, TC-1 together
with the identity $\lambda_{ji}=(\|v\|^2+\langle C^{ji},T\rangle)/4$ exactly for a random rational $v$).
(2) A pure two-qubit vector is $(U\otimes V)(\cos\varphi|00\rangle+\sin\varphi|11\rangle)$ with $0\le\varphi\le\pi/4$; its
correlation matrix is $\operatorname{diag}(\sin2\varphi,-\sin2\varphi,1)$ in the frames rotated by $U,V$, and the concurrence
$\sin2\varphi=|\sum_jv_j^2|=\sqrt\tau$ (the coefficient matrix of $Mv$ has determinant $\sum_jv_j^2/2$).  The local Bloch
vectors are $\cos2\varphi\,r$ and $\cos2\varphi\,s$, of length $\sqrt{1-\tau}$.
(3) $m_\Pi$ is affine in $A$.  The first formula is \eqref{eq:fqe-critical-bures-orbit}.  The second is the overlap of
$\psi$ with a product ray with Bloch vectors $\varepsilon/\sqrt3,\eta/\sqrt3$:
$\langle\psi|\tfrac12(1+\hat\varepsilon\cdot\sigma)\otimes\tfrac12(1+\hat\eta\cdot\sigma)|\psi\rangle
=\tfrac14(1+\hat\varepsilon\cdot r_A+\hat\eta\cdot r_B+\hat\varepsilon^TT\hat\eta)$.
TC-2 and TC-3 check the affine representation of $A$, the equality of the two maxima for 200 random states, and the singular
values $(1,\kappa,\kappa)$, $\det T=-\kappa^2$.
\end{proof}
Thus $\tau=|w|^2$, and a critical state is determined (up to the discrete choices) by $(r,s,w)$.  At $\tau=0$ one recovers the
auxiliary Bloch vectors of \eqref{eq:fqe-null-critical-bloch}.  In these variables
\[
 \begin{gathered}
 \delta=\mathrm{EMAX}-\tfrac13\sum_{j,i}\operatorname{ent}\bigl(\tfrac{1+t_{ji}}4\bigr),\qquad
 d^2=2-2\sqrt{\max_\Pi\Phi_\Pi},\\
 G:=\delta-C_*d^2=\max_\Pi G_\Pi,\qquad G_\Pi=\delta-C_*(2-2\sqrt{\Phi_\Pi}),
 \end{gathered}
\]
with $\mathrm{EMAX}=1+\frac12\log_23$ and $\operatorname{ent}(x)=-x\log_2x$.  Each $G_\Pi$ is analytic wherever all $\lambda_{ji}>0$ and
$\det m_\Pi>0$, and $\delta$ is a convex function of $T$.

\subsubsection*{The chart}
For $\zeta_1,\zeta_2,\zeta_1',\zeta_2'\in[0,1]$ and a sign $\mathrm{sg}=\pm1$ put
\begin{align*}
 r&=\tfrac1N(1,\zeta_1,\zeta_2),\ x_a=\tfrac1n(\zeta_1,-1,0),\ y_a=r\times x_a,\\
 s&=\tfrac1{N'}(\mathrm{sg},\zeta_1',\zeta_2'),\ x_b=\tfrac1{n'}(\zeta_1',-\mathrm{sg},0),\ y_b=s\times x_b,
\end{align*}
$N=\sqrt{1+\zeta_1^2+\zeta_2^2}$, $n=\sqrt{1+\zeta_1^2}$ (and primed), and for $w\in\mathbb R^2$, $|w|\le1$,
\begin{equation}\label{eq:fqe-critical-chart}
 T(y)=rs^T+w_1(x_ax_b^T-y_ay_b^T)+w_2(x_ay_b^T+y_ax_b^T),\qquad y=(\zeta_1,\zeta_2,\zeta_1',\zeta_2',w_1,w_2).
\end{equation}
\begin{lemma}[Reduction to the chart; analytic]\label{lem:fqe-critical-chart}
$\delta$ and $d$ are invariant under $T\mapsto S_aTS_b^T$, where $S_a=D_a\Pi_a$, $S_b=D_b\Pi_b$ with
$\Pi_a,\Pi_b$ cyclic permutation matrices and $D_a,D_b$ diagonal sign matrices, $\det D_a\det D_b=1$ (an even total number of
sign changes), and under the transposition $T\mapsto T^T$.  Every critical state is equivalent under such maps to one
with $T=T(y)$, $y$ in the chart $[0,1]^4\times\{|w|\le1\}$ and $\mathrm{sg}=\pm1$, and then $\tau=|w|^2$.
In the chart the transposition is $(\zeta_1,\zeta_2,\zeta_1',\zeta_2',w_1,w_2)\mapsto(\zeta_1',\zeta_2',\zeta_1,\zeta_2,w_1,\mathrm{sg}\,w_2)$, so it suffices to treat $\zeta_1\le\zeta_1'$.
\end{lemma}
\begin{proof}
The $t_{ji}$ are permuted: this is a linear identity in $T$, and holds on the rank-one matrices $rs^T$ by
Lemma~\ref{lem:fqe-p1-sym}, which span all matrices.  The set of $32$ cube pairs is invariant (PM-2: the parity
$\prod\varepsilon_j\eta_j$ is preserved because $\det D_a\det D_b=1$), the local vectors transform with $S_a,S_b$ up to a
common sign which is absorbed by $(\varepsilon,\eta)\mapsto-(\varepsilon,\eta)$, and $\varepsilon^TT\eta$ is invariant;
so $\Phi$ is invariant by Lemma~\ref{lem:fqe-critical-correlation}(3).  Given $T$, take $r,s$ as in
Lemma~\ref{lem:fqe-critical-correlation}(2); cyclic shifts of $r$ alone and of $s$ alone make $|r_1|,|s_1|$ maximal, and
flips of an even number of the six coordinate signs make $r_1,r_2,r_3,s_2,s_3\ge0$ (the sign of $s_1$ is then determined and
is $\mathrm{sg}$).  The sign patterns of $r$ and of $s$ may have $\det D_a=\det D_b=-1$; then $D_a=-D_a'$, $D_b=-D_b'$ with
$\det D_a'=\det D_b'=1$ and $D_aTD_b^T=D_a'TD_b'^T$, i.e.\ the common sign flip of $(r,s)$, which leaves $rs^T$ fixed, is absorbed.
$(r,s)$ is determined by $T$ up to this common sign when $\tau<1$; at $\tau=1$ the top singular value is degenerate and $r,s$ are
not unique, but $\tau=1$ lies in the analytic range and the chart is used only for $\tau\le0.37$.  Then $r=(1,\zeta_1,\zeta_2)/N$, $s=(\mathrm{sg},\zeta_1',\zeta_2')/N'$ with $\zeta_1,\zeta_2,\zeta_1',\zeta_2'\in[0,1]$, and $T-rs^T$ lies in the
span of the two matrices of \eqref{eq:fqe-critical-chart} by Lemma~\ref{lem:fqe-critical-correlation}(2), with $|w|=\kappa$.
Transposition permutes the twelve $t_{ji}$ (again a linear identity, on rank-one matrices Lemma~\ref{lem:fqe-p1-sym} gives the swap
$r\leftrightarrow s$) and exchanges the roles of the two cube vertices in $\Phi$. Its action on the chart coordinates follows by
direct substitution: $T(y)^T=DT(y')D$, with $y'$ the image coordinates displayed in the lemma and $D=I$ for $\mathrm{sg}=1$,
$D=\operatorname{diag}(1,-1,-1)$ for $\mathrm{sg}=-1$, an admissible sign change (TC-4 also checks it numerically).
\end{proof}

\subsubsection*{Large $\tau$}
\begin{lemma}[Analytic range; TC-6]\label{lem:fqe-critical-large-tau}
If $\tau\ge0.36$ then $\delta\ge C_*d^2$, strictly.
\end{lemma}
\begin{proof}
Put $k=(\ln3-1)/(2\ln2)$.  By \eqref{eq:fqe-critical-entropy} and \eqref{eq:fqe-critical-fidelity},
$\delta\ge k(1-\Phi+5\tau/2)$, and $\Phi\ge f_{\rm cov}(\tau)$ by \eqref{eq:fqe-critical-fidelity-floor}.  With
$x=\sqrt{\Phi}$ the function $k(1-x^2)+5k\tau/2-2C_*(1-x)$ is concave in $x$, hence on $[\sqrt{f_{\rm cov}},1]$ it
is at least its smaller endpoint value; at $x=1$ it is $5k\tau/2>0$.  At $x=\sqrt{f_{\rm cov}(\tau)}$, writing
$\tau=1-u^2$, $f=f_{\rm cov}=\frac13+\frac u{2\sqrt3}$, it is
$\psi(u)=k(1-f)+\frac52k(1-u^2)-2C_*(1-\sqrt f)$, which is concave on $0\le u\le4/5$: $k(1-f)$ is affine,
$-\frac52ku^2$ is concave, and $\sqrt f$ is the square root of a positive affine function, hence concave, with coefficient
$2C_*>0$.  So $\psi$ is positive on $[0,4/5]$ once it is positive at the two endpoints, which TC-6 encloses in ball
arithmetic ($\psi(0)>0.064$, $\psi(4/5)>1.7\cdot10^{-4}$).
\end{proof}

\subsubsection*{The local cone lemma}
At $w=0$ the equality points of Theorem~\ref{thm:fqe-p1-sharp} lie in the chart at the eight points
$(\zeta_1,\zeta_2,\zeta_1',\zeta_2')\in\{(0,0,0,1),(0,0,1,0),(0,1,0,0),(1,0,0,0)\}$, one set for each sign $\mathrm{sg}$.
Let $z_i=|y_i-y^{\rm eq}_i|$ for the four chart coordinates and $z_{5,6}=w_{1,2}$; the box $\mathcal B_\rho$ is
$|z|_\infty\le\rho$, $\rho=1/256$.
\begin{lemma}[Local cone; LCC]\label{lem:fqe-critical-local-cone}
At each equality point exactly four planes $\Pi$ are tied: $G_\Pi=0$ and the others are below $-0.05$.  Their gradients (in $z$)
sum to $0$ and have rank three.  In $\mathcal B_\rho$ one has $G>0$ except at the point itself.
\end{lemma}
\begin{proof}
By Taylor's formula with the Hessian enclosed over the box, $G_\Pi(t\omega)=t\,g_\Pi\!\cdot\!\omega+\frac{t^2}2\omega^T H_\Pi(\xi)\omega$
for $\omega$ on the unit sphere of $|\cdot|_\infty$ (in the orthant of the four chart coordinates) and $0<t\le\rho$, with
$H_\Pi(\xi)$ in an interval matrix $[H_\Pi]$, and $G\ge\sum_\Pi\mu_\Pi G_\Pi$ for weights $\mu_\Pi\ge0$ summing to one on the tied planes.
Write $\mu=\frac14+\nu$.  Because $\sum g_\Pi=0$, the first-order term of the average depends on $\nu$ only.  The tie and this identity are exact, and
are proved symbolically (EX, no floating point or enclosures) at all eight points: exactly four planes have $\Phi_\Pi=\Phi_*$, where
$4\Phi_*=(1+1/\sqrt3)(1+\sqrt{2/3})$ and all other $\Phi_\Pi$ are smaller; the twelve probabilities are $(2\pm\sqrt2)/8$ and $\frac14$, so
$\delta=\frac12\log_23-\frac13(1+2H_{\rm b}(p))$ and $\delta=C_*(2-2\sqrt{\Phi_*})$, i.e.\ $G_\Pi=0$ at the point for the four tied planes, which
is the definition of $C_*$ in \eqref{eq:fqe-p1-cstar}; and $\sum_{\Pi\ \rm tied}\partial_kG_\Pi=0$ for all six $k$, as an identity in
$C_*,\ln p,\ln(1-p),\ln2$.  The reason is the sign symmetry
$z\mapsto(s_1z_1,s_2z_2,s_1s_2z_3,z_4,s_2z_5,s_1z_6)$ at $(0,0,0,1)$, which EX also checks exactly: $T(\sigma y)=DT(y)D'$ with
$D=\operatorname{diag}(1,s_1,s_2)$, $D'=\operatorname{diag}(1,s_1s_2,1)$ permutes the twelve $t_{ji}$ inside each cut and the sixteen
planes, so $G_\Pi\circ\sigma=G_{\pi(\Pi)}$ and only the $z_4$ component of $\sum_\Pi g_\Pi$ can survive; the $z_4$ component vanishes in EX
(it is the derivative of the wall function $G_w$ at $\pi/4$, zero by the reflection symmetry of Lemma~\ref{lem:fqe-p1-wall}).  By the sign symmetry the tied
gradients have the sign form $(s_1a,s_2a,-s_1s_2c,0,-s_2d,-s_1d)$ of the remark below, so their rank is exactly three
(the characters $s_1,s_2,s_1s_2$ carry independent vectors, as the enclosures give $a,c\ne0$); the bound $-0.05$ for the other planes
and the Hessian enclosures are ball arithmetic (LCC).  A piece $\Omega$ of the sphere is
discharged if, for some $\nu$, $\ell=\sum\nu_\Pi g_\Pi\!\cdot\!\omega\ge0$ and $\ell+\rho q>0$ (or $q>0$) with
$q=\frac12\omega^T[H_\mu]\omega$, evaluated in ball arithmetic. This suffices: on $\Omega$ the average is at least
$t\ell+t^2q=t(\ell+tq)$, and $\ell+tq$ is positive for $0<t\le\rho$, being affine in $t$ with $\ell\ge0$ at $t=0$ and
$\ell+\rho q>0$ at $t=\rho$ (or increasing if $q>0$).  LCC finds such weights on $9012$ cells for the eight points.
\end{proof}

\subsubsection*{The covering}
\begin{proposition}[Interval covering; CCV]\label{prop:fqe-critical-covering}
Let $\mathcal R$ be the part of the chart with $|w|^2\le37/100$, both signs $\mathrm{sg}$.  Every box of the bisection tree is
discharged by one of: (G) a certified positive lower bound of $\delta-C_*d^2$; (HS) $d^2\le1/400$ (then
$\delta\ge\frac3{10}d^2$, strictly greater than $C_*d^2$ when $d>0$, by Theorem~\ref{thm:fqe-orbit-stability}); (core) it lies in a box $\mathcal B_\rho$ of
Lemma~\ref{lem:fqe-critical-local-cone}.
\end{proposition}
\begin{proof}
The starting closed boxes cover $[0,1]^4\times[-5/8,5/8]^2$ in each sign sector, hence the required disc.
A split replaces one box by two closed children whose union is the parent; their common boundary is retained.
Boxes crossing the disc or ordering boundary are retained unless they are wholly outside the domain reduced by symmetry.
Thus complete coverage means that every retained branch ends in one of the three certified alternatives in the statement.
A failed or inconclusive sign test only triggers subdivision; it never certifies a box.

There is no need to choose a unique maximising fidelity plane. Since $G=\max_\Pi G_\Pi$,
any convex combination $G_\mu$ is a lower bound for $G$, including on a tie.
Near equality, the exact contacts and gradient cancellation of the local cone lemma replace a strict sign test at the apex.

On a box with centre $c$ and half-widths $h$ (in $y$) use the weights $\mu$ uniform on the best one to four planes.
(i) $\delta$ is convex in the twelve numbers $t_{ji}$, hence $\delta(y)\ge\delta(c)+\sum_iD_i(t_i(y)-t_i(c))$ with
$D_i=(1+\ln\lambda_i(c))/(12\ln2)$, valid for every point of the box since every chart point is a state, so that $t_i(y)\ge-1$
(the tangent of a convex function holds on its whole domain); a ray whose probability at $c$ is not certified positive is left
out of the tangent and bounded by the maximum of $-x\log_2x$ over its probability box ($-x\log_2x$ increases on $[0,1/e]$).
(ii) $\Phi_\Pi$ is concave in the three entries $m$ of $m_\Pi$, so $d_\Pi^2=2-2\sqrt{\Phi_\Pi}$ is convex there:
$d_\Pi^2(m)\le d_\Pi^2(m_c)+a_\Pi\!\cdot\!(m-m_c)+B$, where the Bregman remainder $B$ is convex in $m-m_c$ and so at most its maximum at the
corners of the box $|m-m_c|\le r$; $r$ is a bound for $|m(y)-m(c)|$.
(iii) $t(y)-t(c)$ and $m(y)-m(c)$ are expanded in $y-c$ to second order, with the exact gradient at $c$ and the interval second
derivatives of the nine entries $T_{\alpha\beta}(y)$ over the box (closed forms, ball arithmetic).
With $g_k=\partial_kG_\mu(c)$ this gives $G\ge G_\mu(c)-\sum_k|g_k|h_k-R$, where $R$ is the sum of the second-order terms of (iii)
inserted into the tangent of (i) and, with factor $C_*$, into the linear term of (ii), and of $C_*\max B$, the first-order terms being exact.
Bisection of the widest side (the floating-point routine \texttt{critical\_float.py} only orders the attempts and certifies nothing)
discharges the chart region $\zeta_1\le\zeta_1'$, $|w|^2\le37/100$ with $208953$ boxes (G) in the sector $\mathrm{sg}=+1$, and
$246044$ (G) and $436$ (HS) for $\mathrm{sg}=-1$, besides $42$ boxes inside the local cone boxes in each sector;
the ordering comparison $\zeta_1>\zeta_1'$ is exact on the dyadic endpoints, while the outside-disc screen has a conservative
positive margin exceeding its floating-point roundoff, so it discards only boxes strictly outside $|w|^2\le37/100$.
The orbit point of the chart ($d=0$, for
$\mathrm{sg}=-1$) is covered by (HS).
\end{proof}

\begin{proof}[Proof of Theorem~\ref{thm:fqe-p1-sharp} for critical states]
If $\tau\ge0.36$, apply Lemma~\ref{lem:fqe-critical-large-tau}.  Otherwise Lemma~\ref{lem:fqe-critical-chart} reduces to
the chart with $|w|^2\le37/100$, which Proposition~\ref{prop:fqe-critical-covering} and Lemma~\ref{lem:fqe-critical-local-cone}
cover: $G>0$, or $d>0$ and $\delta\ge\frac3{10}d^2>C_*d^2$, or $d=0$, except at the equality points $y^{\rm eq}$ where $G=0$ by
\eqref{eq:fqe-p1-cstar}.  The equality points have $w=0$ and are the $144$ points of Theorem~\ref{thm:fqe-p1-sharp}.
\end{proof}

\begin{remark}[Structure at the equality points; numerical picture]\label{rem:fqe-critical-structure}
In the coordinates $z=(\zeta_1,\zeta_2,\zeta_1',1-\zeta_2',w_1,w_2)$ at the equality point with $(\zeta_1,\zeta_2,\zeta_1',\zeta_2')=(0,0,0,1)$, the four tied
plane functions have the gradients $g=(s_1a,\,s_2a,\,-s_1s_2c,\,0,\,-s_2d,\,-s_1d)$, $s_1,s_2=\pm1$, with
$a=0.059039$, $c=0.036251$, $d=0.018765$ (ball enclosures, LCC).  They sum to zero and have rank three, with singular values
$0.1239,0.1239,0.0725$: $G$ has a conical minimum in three directions and is smooth in the other three.  The average of the
four tied Hessians is positive definite on the three-dimensional kernel of the gradients, with eigenvalues
$0.0725$ (the wall direction $\theta$) and $0.7285,0.7285$ (two combinations of the chart coordinates with the $w$-plane).
In particular the dependence on $w$ is first order, in the shape of a square $(\pm1,\pm1)$, so that the loss of
fidelity under moving off $\tau=0$ is compensated at first order by the cone and costs nothing along the smooth directions.
Numerically (not a proof; numerical evidence, NV-1) the minimum of $\delta-C_*d^2$ over critical states with fixed
$\tau$ is about $0.31\,\tau$ for every $\tau\in(0,1]$ (about $0.30\tau$ at $\tau=0.4$ and $0.5$), attained near the equality points for small $\tau$ and,
at $\tau=1$, at a real state with $d^2=2-\sqrt2$.
\end{remark}

\begin{remark}[Non-critical states; numerical evidence only]\label{rem:fqe-critical-noncritical}
The theorem is for critical states; whether $C_*$ is optimal over all states is open.  Evidence (not a proof):
at the equality state, random positive local filters of size $\varepsilon\le0.2$ (the twelve non-critical directions of the
$18$-dimensional LU quotient; $12$ directions per size) give $\delta-C_*d^2\ge3\varepsilon^2$ (NV-2); on states obtained by filtering critical states within $0.4$ of the equality state
(normal-form coordinates) with filters of size $0.05$ to $0.4$ the minimum of $\delta/d^2$ is $0.27$.
These sampled tests found no value below $C_*$. Reduction of a non-critical state to its critical representative does not preserve
$\delta$ (by Section~\ref{sec:fqe-renyi-filter-limit}, one filter on one qubit of a critical state cannot raise any
balanced-cut entropy, but filters on two qubits can raise the mean R\'enyi entropy of order $3/2$), and the general coefficient $k(6+\sqrt6)/12\approx0.05$ is a factor $3.8$ below
$C_*$, so a proof needs a quantitative treatment of the filter directions: near the critical locus a Taylor bound in the filter
variables over this six-parameter chart (a covering of the same size), far from it a stronger distance--entropy comparison.
\end{remark}
